% hybrid_extractor の動作確認用サンプル。
% IR版3（common/README.md）の項目がひと通り出るように書いている。
%
% PDFにする（ai/hybrid_extractor/sample_data/ で実行）:
%   uplatex sample.tex
%   uplatex sample.tex     ← 2回目。表・図の番号や参照を確定させる
%   dvipdfmx sample.dvi    → sample.pdf
%
% どの部分が、IRのどの項目を確かめるためのものかは、各所のコメントに書いた。

\documentclass[uplatex,a4paper,11pt]{jsarticle}
\usepackage[dvipdfmx]{graphicx}
\usepackage{amsmath,amssymb}
\usepackage{booktabs}  % 縦線の無い表（\toprule など）
\usepackage{array}     % 表の列の指定を増やす

% label: title / author / date
\title{ハイブリッド抽出器のためのサンプル文書}
\author{谷本 龍信}
\date{2026年11月8日}

\begin{document}
\maketitle

% label: heading（level 1）
\section{はじめに}

% label: paragraph（2行以上になる、普通の段落）
この文書は、PDFから文字を取り出してブロックに分ける抽出器と、そのブロックの種類を当てる分類器を、
開発しながら確かめるためのサンプルである。見出し・段落・箇条書き・数式・表・図・コードを、
それぞれ少なくとも1つずつ含むように書いている。

% label: paragraph ＋ runs（太字・下線。日本語の書体には斜体が無いので、斜体は英語で入れる）
段落の中には、\textbf{太字の語}や\underline{下線を引いた語}が混ざることがある。
英語では \textit{italic words} も使われる。抽出器は、こうした装飾が段落のどの文字に付いているかを
runs として出す。

% label: heading（level 2）
\subsection{箇条書き}

% label: list（記号の付いた箇条書きと、番号の付いた箇条書き）
抽出器が確かめるべきことは、次のとおりである。
\begin{itemize}
  \item 文字の大きさとフォントを取り出せること
  \item 行を正しくつないで段落にできること。この項目はわざと長くして、2行に折り返されるようにしている
  \item 表と図の領域を見つけられること
\end{itemize}

作業の順番は次のとおりである。
\begin{enumerate}
  \item PDFを開いて、文字を取り出す
  \item 文字をブロックにまとめる
  \item ブロックの種類を決める
\end{enumerate}

% label: heading（level 3）
\subsubsection{文中の数式}

% label: paragraph（文中の数式。段落の文字のままになる）
文中に $E = mc^2$ や $a_i + b_i$ のような数式が入ることもある。これは別行立ての数式とは分けて扱う。

\section{数式とコード}

% label: math（別行立ての数式、2行）
別行立ての数式の例を示す。
\begin{align*}
  f(x) &= \sum_{i=1}^{n} a_i x^i \\
  \int_0^1 x^2 \, dx &= \frac{1}{3}
\end{align*}

% label: code（字下げを含むコード。text には行頭の空白を残す）
コードの例を示す。
\begin{verbatim}
def main():
    for page in doc:
        if page.number == 0:
            print(page.get_text())
\end{verbatim}

\section{表と図}

% kind: table ＋ cells / aligns / col_rules / row_rules
% 列の揃えは 左・中央・右。縦線は全部、横線は上・見出しの下・下だけ（行の間には無い）
表\ref{tab:grid}は、罫線のある表の例である。
\begin{table}[htbp]
  \centering
  \caption{罫線のある表}  % label: caption
  \label{tab:grid}
  \begin{tabular}{|l|c|r|}
    \hline
    項目 & 担当 & 工数 \\
    \hline
    抽出器 & 谷本 & 15時間 \\
    分類器 & 谷本 & 20時間 \\
    評価 & 加茂 & 8時間 \\
    \hline
  \end{tabular}
\end{table}

% kind: table（縦線の無い表。booktabs）＋ col_widths（2列目はセルの中で折り返す幅指定の列）
表\ref{tab:booktabs}は、縦線の無い表と、セルの中で文字を折り返す列の例である。
\begin{table}[htbp]
  \centering
  \caption{縦線の無い表}
  \label{tab:booktabs}
  \begin{tabular}{lp{5cm}}
    \toprule
    手法 & 説明 \\
    \midrule
    ルールベース & 文字の大きさや太字などの決まりで、ブロックの種類を決める。学習はいらない \\
    自作AI & 正解の付いた文書から学習して、ブロックの種類を当てる \\
    \bottomrule
  \end{tabular}
\end{table}

% kind: figure ＋ caption（画像ファイルを使わず、枠線だけの図にしている）
図\ref{fig:box}は、図の例である。
\begin{figure}[htbp]
  \centering
  \fbox{\rule{0pt}{3cm}\hspace{0.5\linewidth}}
  \caption{枠線だけの図}
  \label{fig:box}
\end{figure}

\section{おわりに}

% label: paragraph（2ページ目に続く文章。ページをまたいだときのブロックを確かめる）
この文書を使えば、抽出器を書き換えるたびに同じ入力で結果を比べられる。
ランダムに作る評価用のデータセットとは違い、中身が毎回同じなので、どこが変わったのかを目で確かめやすい。
評価用のデータセットで点数を測り、このサンプルで原因を探す、という使い分けをする。
ページ番号は、jsarticle の既定のとおり各ページの下の中央に出る。
% ↑ ページ番号は label: other

\end{document}
